\section{Circuit Complexity --- UHalt}
%%% generated by the blueprint pipeline from TCSlib.Complexity.CircuitComplexity.UHalt
\subsection*{Circuit Complexity --- UHalt: Introduction}

Declarations of the Lean module \texttt{TCSlib.Complexity.CircuitComplexity.UHalt}.

\subsection*{Circuit Complexity --- UHalt: Results}

\begin{definition}[Halting set for program codes]
\lean{Nat.Partrec.Code.haltingSet}
\label{Nat.Partrec.Code.haltingSet}
\leanok
The halting set is the set of natural numbers $n$ such that the partial recursive
program with code given by the first component of the standard pairing decomposition of
$n$ is defined on input given by the second component.
\end{definition}

\begin{theorem}[Undecidability of the halting set]
\lean{Nat.Partrec.Code.not_computablePred_mem_haltingSet}
\label{Nat.Partrec.Code.not_computablePred_mem_haltingSet}
\leanok
\uses{Nat.Partrec.Code.haltingSet}
\difficulty{5}
The set of natural numbers whose encoded program halts on its encoded input is not a
computable predicate.
\end{theorem}

\begin{theorem}[Noncomputability of unary languages]
\lean{Language.not_computablePred_mem_unary}
\label{Language.not_computablePred_mem_unary}
\leanok
\uses{Language.replicate_mem_unary_iff, Language.unary}
\difficulty{4}
Let $S$ be a set of natural numbers whose membership predicate is not computable. The
language of all-true Boolean lists whose lengths belong to $S$ also has no computable
membership predicate.
\end{theorem}

\begin{definition}[Unary halting language]
\lean{Language.uhalt}
\label{Language.uhalt}
\leanok
\uses{Language.unary, Nat.Partrec.Code.haltingSet}
The unary halting language is the language over the Boolean alphabet consisting of the
unary encodings of natural numbers in the halting set; the unary encoding of $n$ is the
word of $n$ copies of the symbol $1$.
\end{definition}

\begin{theorem}[Unary halting language has polynomial layered circuits]
\lean{Language.uhalt_inLayeredPPoly}
\label{Language.uhalt_inLayeredPPoly}
\leanok
\uses{Language.InLayeredPPoly, Language.uhalt, Language.unary_inLayeredPPoly}
\difficulty{1}
Let $H$ be the set of natural numbers $m$ for which the machine with index given by the
first component of the Cantor pairing decomposition of $m$ halts on input given by its
second component. The language of all-true Boolean lists whose lengths belong to $H$ is
recognized by a family of layered circuits over Boolean values, with every layer finite
and using only identity, NOT, and unbounded-fan-in AND gates. There are natural numbers
$a$ and $k$ such that, for every input length $n$, the number of noninput circuit nodes
is at most $a(n+1)^k$.
\end{theorem}

\begin{theorem}[Noncomputability of the unary halting language]
\lean{Language.not_computablePred_mem_uhalt}
\label{Language.not_computablePred_mem_uhalt}
\leanok
\uses{Language.not_computablePred_mem_unary, Language.uhalt, Nat.Partrec.Code.not_computablePred_mem_haltingSet}
\difficulty{1}
Let $H$ be the set of natural numbers $m$ for which the machine with index given by the
first component of the Cantor pairing decomposition of $m$ halts on input given by its
second component. The language consisting of the all-true Boolean lists whose lengths
belong to $H$ has no computable membership predicate.
\end{theorem}

\begin{theorem}[A noncomputable language in layered polynomial size]
\lean{Language.exists_le_allOnes_inLayeredPPoly_not_computablePred}
\label{Language.exists_le_allOnes_inLayeredPPoly_not_computablePred}
\leanok
\uses{Language.InLayeredPPoly, Language.allOnes, Language.not_computablePred_mem_uhalt, Language.uhalt, Language.uhalt_inLayeredPPoly, Language.unary_le_allOnes}
\difficulty{1}
There is a language $L$ of Boolean strings contained in the language of all-ones strings
such that $L$ is recognized by a family of layered circuits using only the specified
standard gate operations and having size bounded by a polynomial in the input length,
but membership in $L$ is not a computable predicate.
\end{theorem}

\begin{theorem}[Polynomial circuit complexity of the halting language]
\lean{Language.uhalt_inPPoly}
\label{Language.uhalt_inPPoly}
\leanok
\uses{Language.InLayeredPPoly.inPPoly, Language.InPPoly, Language.uhalt, Language.uhalt_inLayeredPPoly}
\difficulty{1}
The universal halting language belongs to polynomial size: there are natural numbers $a$
and $k$ and a family of Boolean directed acyclic circuits, one for each input length
$n$, that computes this language, has well-formed gates of fan-in at most two, and has
size at most $a(n+1)^k$ for every $n$.
\end{theorem}

\begin{theorem}[A noncomputable language in polynomial circuit size]
\lean{Language.exists_le_allOnes_inPPoly_not_computablePred}
\label{Language.exists_le_allOnes_inPPoly_not_computablePred}
\leanok
\uses{Language.InPPoly, Language.allOnes, Language.not_computablePred_mem_uhalt, Language.uhalt, Language.uhalt_inPPoly, Language.unary_le_allOnes}
\difficulty{1}
There is a language $L$ of Boolean strings contained in the language of all-ones strings
such that $L$ is recognized by a family of well-formed Boolean circuits with at most two
inputs per gate and size bounded by a polynomial in the input length, but membership in
$L$ is not a computable predicate.
\end{theorem}
